\documentclass[aspectratio=169,10pt]{beamer}
\newcommand{\reporoot}{../..}
% Shared preamble of the main deck and the topic decks.
% Define \reporoot (the path from the .tex file to the repository root)
% before \input-ing this file.
\usetheme[progressbar=frametitle,numbering=fraction,sectionpage=progressbar]{metropolis}

\usepackage{amsmath,amssymb}
\usepackage{booktabs}
\usepackage{array}
\newcolumntype{L}[1]{>{\raggedright\arraybackslash}p{#1}}
\usepackage{graphicx}
\usepackage{tikz}
\usetikzlibrary{arrows.meta,positioning,decorations.pathreplacing}
\usepackage[authoryear,round]{natbib}
\usepackage{xcolor}

\graphicspath{{\reporoot/figures/}}

% Every simulated number on the slides is produced by experiments/*.py and read from here.
\input{\reporoot/results/e01_estimators.tex}
\input{\reporoot/results/e02_scaling.tex}
\input{\reporoot/results/e03_var_backtests.tex}
\input{\reporoot/results/e04_es_backtest.tex}
\input{\reporoot/results/e05_pla.tex}
\input{\reporoot/results/e06_estimation_risk.tex}
\input{\reporoot/results/e07_model_uncertainty.tex}
\input{\reporoot/results/e08_asrf.tex}
\input{\reporoot/results/e09_copula_is.tex}
\input{\reporoot/results/e10_calibration_tests.tex}
\input{\reporoot/results/e11_discrimination.tex}
\input{\reporoot/results/e12_exposure.tex}
\input{\reporoot/results/e13_cva_wwr.tex}
\input{\reporoot/results/e14_nested.tex}
\input{\reporoot/results/e15_variance_reduction.tex}
\input{\reporoot/results/e16_lda.tex}
\input{\reporoot/results/e17_evt.tex}
\input{\reporoot/results/e18_reverse_stress.tex}
\input{\reporoot/results/e19_regimes.tex}
\input{\reporoot/results/gallery.tex}

\newcommand{\E}{\mathbb{E}}
\newcommand{\Var}{\operatorname{Var}}
\newcommand{\VaR}{\mathrm{VaR}}
\newcommand{\ES}{\mathrm{ES}}
\newcommand{\dd}{\mathrm{d}}
\definecolor{accent}{HTML}{0C5DA5}
\definecolor{warn}{HTML}{C0392B}
\newcommand{\hl}[1]{\textcolor{accent}{\textbf{#1}}}

\setbeamerfont{footnote}{size=\tiny}
\setbeamerfont{caption}{size=\scriptsize}

\author{Yaohui Shu}
\institute{Stony Brook University}
\date{November 2026}

% Topic decks: a references frame that lists only the works cited in the deck.
\newcommand{\deckreferences}{%
  \begin{frame}[allowframebreaks]{References}
    \tiny
    \renewcommand{\bibsection}{}
    \bibliographystyle{plainnat}
    \bibliography{\reporoot/slides/refs}
  \end{frame}}

% Topic decks: \decknumber{n} puts the deck number on the title page.
\newcommand{\decknumber}[1]{\institute{Stony Brook University\hfill{\scriptsize\color{gray}Deep dive #1 of 9}}}


\title{Backtesting VaR and Expected Shortfall}
\subtitle{Exception counts, independence, ES tests and the FRTB P\&L attribution test}
\decknumber{2}

\begin{document}

\maketitle

\begin{frame}{The question}
  \small
  \begin{columns}[T]
    \begin{column}{0.55\textwidth}
      A risk model forecasts a tail every day. A year later we have 250 outcomes. \hl{Can we tell a good model from a bad one, and how often do we get it wrong?}
      \begin{itemize}
        \item Regulators turn backtest outcomes into capital: more exceptions, a higher multiplier.
        \item Validation teams need tests with known size (false alarms) and known power (missed problems).
      \end{itemize}
      Simulation gives both, because we control the truth.
    \end{column}
    \begin{column}{0.42\textwidth}
      \begin{alertblock}{In this deck}
        \small
        \begin{enumerate}
          \item Exception tests and the Basel traffic light.
          \item What 250 days can detect.
          \item Backtesting ES.
          \item The FRTB P\&L attribution test.
        \end{enumerate}
      \end{alertblock}
    \end{column}
  \end{columns}
\end{frame}

\begin{frame}{Backtesting as a random experiment}
  \centering
  \includegraphics[width=0.95\textwidth,height=0.54\textheight,keepaspectratio]{gallery/g02_backtesting.png}\par\vspace{0.4em}
  \footnotesize\raggedright
  (a) Exception counts in 250 days for a correct 99\% VaR (true exception rate $p = 1\%$) and for a model with $p = 2\%$, over the Basel zones. (b) Rejection probability of the exact one-sided binomial test of $p = 1\%$ at 5\% size, against the true rate $p$ and the backtest length $n$ (days). Against $p = 2\%$ it rejects \GalTwoPowTwoFifty\% of the time at 250 days; 80\% power (red) takes about \GalTwoNEighty{} days.
\end{frame}

\begin{frame}{Exception-based tests}
  \small
  Hit $I_t = \mathbf{1}\{L_t > \VaR_{t}\}$. Under a correct VaR at level $\alpha$, the $I_t$ are i.i.d.\ Bernoulli($p$), $p = 1 - \alpha$. With $x$ exceptions in $n$ days:
  \begin{columns}[T]
    \begin{column}{0.5\textwidth}
      \textbf{Proportion of failures} \citep{kupiec1995techniques}:
      \[ LR_{pof} = -2\ln\frac{(1-p)^{n-x}p^x}{(1-\hat p)^{n-x}\hat p^{\,x}} \sim \chi^2_1,\quad \hat p = x/n. \]
      \textbf{Independence} \citep{christoffersen1998evaluating}: a two-state Markov chain with $\pi_{ij} = P(I_t = j \mid I_{t-1} = i)$ against $\pi_{01} = \pi_{11}$; $LR_{ind} \sim \chi^2_1$, and conditional coverage $LR_{cc} = LR_{pof} + LR_{ind} \sim \chi^2_2$.
    \end{column}
    \begin{column}{0.46\textwidth}
      \textbf{Basel traffic light} \citep{bcbs2023mar32}, 99\% VaR, 250 days, multiplier added to $m_c$:
      \vspace{0.2em}

      \footnotesize
      \begin{tabular}{lll}
        \toprule
        Zone & Exceptions & Multiplier \\
        \midrule
        Green & 0 to 4 & 1.50 \\
        Amber & 5 to 9 & 1.70 to 1.92 \\
        Red & 10 or more & 2.00 \\
        \bottomrule
      \end{tabular}
      \vspace{0.3em}

      \small Exceptions count the larger of actual-P\&L and hypothetical-P\&L exceptions. A desk with more than 12 exceptions at 99\% or 30 at 97.5\% in a year moves to the standardised approach (MAR32.19).
    \end{column}
  \end{columns}
\end{frame}

\begin{frame}{What 250 days can detect}
  \centering
  \includegraphics[width=0.82\textwidth]{fig03_var_backtests.png}
  \footnotesize
  \begin{columns}[T]
    \begin{column}{0.32\textwidth}
      (a) Exact binomial. A correct model lands green $\EThreeGreenOne\%$, amber $\EThreeAmberOne\%$, red $\EThreeRedOne\%$ of the time. A model whose true rate is 2\% is still green $\EThreeGreenTwo\%$ of the time.
    \end{column}
    \begin{column}{0.32\textwidth}
      (b) Correct model, $\chi^2$ critical values, 99\% VaR. At 250 days Kupiec rejects $\EThreeSizeKupShort\%$ and independence $\EThreeSizeIndShort\%$ instead of 5\%; at 10,000 days independence still rejects only $\EThreeSizeIndLong\%$. Use exact or simulated critical values.
    \end{column}
    \begin{column}{0.32\textwidth}
      (c) GARCH-t returns, VaR at the true unconditional quantile: right rate ($\EThreeRate\%$), clustered hits. Clustering makes the count overdispersed, so Kupiec rejects $\EThreePowKupShort\%$ at 250 days. Independence rejects $\EThreePowIndShort\%$, rising to $\EThreePowIndLong\%$ with 10,000 days.
    \end{column}
  \end{columns}
\end{frame}

\begin{frame}{Backtesting expected shortfall}
  \small
  ES is not elicitable alone \citep{gneiting2011making}, but (VaR, ES) is jointly elicitable \citep{fissler2016higher}, and ES tests use both forecasts.
  \vspace{0.3em}
  \begin{columns}[T]
    \begin{column}{0.52\textwidth}
      \textbf{Acerbi and Sz\'ekely (2014), test 2.} With returns $X_t$ (gains positive), forecasts $\VaR_t$, $\ES_t$ at tail probability $p$, and $I_t = \mathbf{1}\{X_t + \VaR_t < 0\}$:
      \[ Z_2 = \sum_{t=1}^{T}\frac{X_t I_t}{T\,p\,\ES_t} + 1, \qquad \E_{H_0}[Z_2] = 0, \]
      negative when risk is underestimated. It weighs each exception by its size relative to the forecast ES \citep{acerbi2014backtesting}.
    \end{column}
    \begin{column}{0.45\textwidth}
      The authors report that critical values are stable across Student $t$ models: about $-0.7$ at 5\% and $-1.8$ at 0.01\% for $T = 250$, so fixed thresholds can serve as a traffic light.
      \vspace{0.3em}

      Regression-based ES backtests that need only the ES forecast are an alternative \citep{bayer2022regression}.
    \end{column}
  \end{columns}
\end{frame}

\begin{frame}{ES test against exception counting}
  \centering
  \includegraphics[width=0.88\textwidth]{fig04_es_backtest.png}
  \footnotesize
  \begin{columns}[T]
    \begin{column}{0.32\textwidth}
      (a) Model and truth Student $t_5$, $T = 250$, 200,000 simulated years. The 5\% critical value $\EFourCritFive$ matches the published $-0.7$. The 0.01\% value, $\EFourCritTiny$ against the published $-1.8$, rests on about 20 simulated years and is imprecise.
    \end{column}
    \begin{column}{0.32\textwidth}
      (b) True scale 1.2 times the model's: $Z_2$ rejects $\EFourPowZTwoK\%$ of the time, the 99\% VaR count test ($x \ge \EFourXc$, size $\EFourSizeVaR\%$) $\EFourPowVaRK\%$.
    \end{column}
    \begin{column}{0.32\textwidth}
      (c) Same variance, heavier tail ($\nu = 3$): $\EFourPowZTwoNu\%$ against $\EFourPowVaRNu\%$. Both tests struggle with tail shape at 250 days.
    \end{column}
  \end{columns}
\end{frame}

\begin{frame}{The FRTB P\&L attribution test}
  \small
  \begin{columns}[T]
    \begin{column}{0.5\textwidth}
      Each trading desk compares, over 250 days, the hypothetical P\&L (HPL) from front-office pricing with the risk-theoretical P\&L (RTPL) from the risk model (MAR32.20 to MAR32.42):
      \begin{itemize}
        \item Spearman rank correlation of HPL and RTPL;
        \item Kolmogorov-Smirnov (KS) distance between their empirical distributions.
      \end{itemize}
      Green if correlation $> 0.80$ and KS $< 0.09$; red if correlation $< 0.70$ or KS $> 0.12$; amber otherwise (capital surcharge). Red desks use the standardised approach.
    \end{column}
    \begin{column}{0.46\textwidth}
      \textbf{Three ways RTPL can differ} (Student $t_4$ factor P\&Ls, 4,000 simulated years per point):
      \begin{itemize}
        \item a missing risk factor with share $f$ of HPL variance;
        \item a scale error, RTPL $= s \times$ HPL;
        \item valuation noise with variance share $g$.
      \end{itemize}
      A pure scale error leaves the rank correlation at one (lowest mean $\EFiveRhoScale$), so only KS can catch it.
    \end{column}
  \end{columns}
\end{frame}

\begin{frame}{PLA zones in simulation}
  \centering
  \includegraphics[width=0.94\textwidth]{fig05_pla.png}
  \footnotesize
  \begin{columns}[T]
    \begin{column}{0.32\textwidth}
      (a) A perfect model is green $\EFiveGreenZero\%$ of the time. Green falls below 50\% once $\EFiveHalfGreenF\%$ of variance is missing.
    \end{column}
    \begin{column}{0.32\textwidth}
      (b) Green falls below 50\% only for $s \ge \EFiveHalfGreenSUp$ or $s \le \EFiveHalfGreenSLo$: with 250 days, the KS threshold tolerates scale errors of 20\% or more before flagging them.
    \end{column}
    \begin{column}{0.32\textwidth}
      (c) Noise: green below 50\% at $g = \EFiveHalfGreenG\%$ of HPL variance.
    \end{column}
  \end{columns}
\end{frame}

\begin{frame}{Discussion and takeaways}
  \small
  \begin{columns}[T]
    \begin{column}{0.5\textwidth}
      \textbf{What the simulations show}
      \begin{itemize}
        \item 250 days of 99\% exceptions cannot separate a 1\% model from a 2\% model reliably.
        \item $\chi^2$ approximations are poor at Basel sample sizes: Kupiec over-rejects and the independence test under-rejects.
        \item Clustered exceptions are nearly invisible to the independence test at 250 days; they show up instead as overdispersed counts.
        \item The ES test adds power against volatility errors; tail shape stays hard to test with one year.
        \item In the PLA test only KS catches scale errors, and only large ones.
      \end{itemize}
    \end{column}
    \begin{column}{0.46\textwidth}
      \textbf{For a validation team}
      \begin{enumerate}
        \item Report size and power of every test at your sample length, by simulation.
        \item Test independence and ES, not only counts.
        \item Diagnose PLA failures by metric: low correlation points to missing factors, high KS to scale or distribution errors.
      \end{enumerate}
      \textbf{Assumptions.} One-day horizon, i.i.d.\ or GARCH truth, no overlap between forecasts.
    \end{column}
  \end{columns}
\end{frame}

\begin{frame}[fragile]{Reproduce}
  \small
  Code: \texttt{backtest.py}. Experiments: \texttt{e03\_var\_backtests.py}, \texttt{e04\_es\_backtest.py}, \texttt{e05\_pla.py}.
\begin{verbatim}
from stochastic_risk import backtest as bt
bt.traffic_light(6)                     # ('amber', 1.76)
bt.kupiec_pof(x=6, n=250, p=0.01)       # (LR, p-value)
bt.christoffersen(hits, p=0.01)         # independence and conditional coverage
bt.acerbi_szekely_z2(returns, var, es, alpha=0.975)
bt.pla_metrics(hpl, rtpl)               # (Spearman, KS, zone)
\end{verbatim}
\end{frame}

\deckreferences
\end{document}
